\begin{proposition}[a global scalar gauge would be trivial]\label{push:src:C.3.2}
Let a smooth circle act on a simply connected smooth manifold $Z$.
If $g:Z\to S^1$ is smooth and invariant under the circle, then
$z\mapsto g(z)\cdot z$ is smoothly isotopic to the identity.
The primitive coordinate $\theta$ detecting $\gamma$ on a regular
fibre of $X_c$ has no continuous extension to $X_c$.
\end{proposition}
\Uses{the hypotheses, notation and definitions stated here; the explicit constructions and calculations in the proof}


\begin{proof}
Since $Z$ is simply connected, $g$ lifts to a smooth real function
$a$ with $g=e^{2\pi ia}$. Invariance implies
$a(\lambda z)-a(z)\in\Z$. This difference is continuous in the
connected circle and is zero at its identity, so it is zero
everywhere. The maps $z\mapsto e^{2\pi ita(z)}\cdot z$ are
diffeomorphisms with inverses obtained by replacing $t$ by $-t$.
They give the isotopy. In the present application a loop representing
$\widehat\gamma$ in a regular torus fibre has degree one under
$\theta$. If $\theta$ extended to the simply connected $S^6$,
that degree would be zero. The contradiction shows why the proposed local scalar twist is not a global scalar-gauge calculation.
\end{proof}

The credited statement and complete re-derivation are in \cref{push:src:C.3.3}.
